\section{Настройка регуляторов}

\subsection{Регулятор тока}

Регулятор тока в настройке не нуждается. Параметры регулятора тока:

\begin{tabular}{|l|l|}
    \hline
    Параметр & Значение \\ \hline
    PID P & 0.5 \\ \hline
    PID I & 0.08 \\ \hline
    PID D & 0 \\ \hline
    WindUp & 0.04 \\ \hline
\end{tabular}

\subsection{Регулятор температуры}

Автоматической настройки регулятора нет. Настройка ручная.
Подготовительные настройки:
\begin{enumerate}
    \item   После включения контроллера установить лимиты на вкладке \textit{Limits} [\ref{sec:Limits}].
            Например для модуля \textbf{TB-127-1.0-2.5, Криотерм} пределы тока будут 1,9 А и минус 1,9 А, соответсвенно,
            а напряжение 15,9 и минус 15,9 В. Пограммное аварийное значение тока отключения установить x2 от рабочего: 4,0 А.

    \item   Проверить правильность включения преобразователя и терморезистора. Для этого переключить контроллер в
            режим стабилизации тока. На вкладке \textit{Current Loop} [\ref{sec:Current}] установить тока 0,1 А и наблюдать рост температуры.
            Возможные проблемы:
            \begin{itemize}
                \item Вместо нагрева идёт охлаждение: поменять местами клеммы подключения преобразователя.
                \item Нет изменения температыры и напряжения на преобразователе близко к ограничению: обрыв в цепи подключения или поломка преобразователя
                \item Температура нестабильна: проверить подключение терморезистораю
            \end{itemize}
    \item   Перевести контроллер в режим Остановки
    \item   Установить скорость нагрева и охлаждения в 1000 К/минуту (только для настройки). Вкладка \textit{Limits} [\ref{sec:Limits}].
    \item   Установить все параметры ПИД регулятора температуры в 0 и передать в контроллер нажатием кнопки Set.
            Для проверки верности установки считать параметры нажатием кнопки Get.
    \item   Перевести регулятор в режим стабилизации тока. Наблюдать установившийся выходно ток около 0 А.
\end{enumerate}

Настройка ПИД регулятора основана на методе Ziegler-Nicholas, на осцилляциях регулятора.

\begin{enumerate}
    \item Задать температуру регулирования на 5-8 градусов больше текущей
    \item Увеличивая пропорциональный коэффициент добиваемся устойчивых колебаний температуры без изменения амплитуды. Начинать можно с коэффициента 0,5.
    \item Записать сигнал на вкладке \textit{Current Loop} [\ref{sec:Current}] кнопкой Start Record. Время записи ~30 секунд.
    \item Записанный сигнал пропускаем через утилиту \textit{pid.py}, она вычислит период и амплитуду колебаний
    \item Полученые значения зносим в таблицу для дальнейших вычислений
    \item Вычисленные коэффициенты ПИД регулятора вносим в настройки при выключенном режиме
\end{enumerate}
